\chapter{More than one key}

One key is one person, one device, one bad day. A
multisig raises the bar only when the signers, the
devices, the places, and the verification habit are
actually independent. A poorly run multisig is theater.

This chapter is the personal and small-team version.
Protocol treasuries, timelocks, pause guardians, and
Safe operator steps live in the Protocol Multisig
companion. Use that book when the wallet is an
institution.

\section{M of N}

You need M signatures out of N keys. 2-of-3 and 3-of-5
are the usual shapes. Avoid N-of-N. Losing one key
should not brick the funds.

The source's numbers (at least three signers, about half
required, more signers above large balances) are
examples from protocol practice. Copy the idea, not the
spreadsheet.

Every signer uses hardware. Dedicated address per
multisig, even if it is only another index on the same
device. Same seed still means one seed incident burns
every seat.

\section{Independence}

Two keys on one laptop is one key. Two devices in one
apartment is one burglary. Two signers who share a
password manager is one phish. Different brands, different
cities, different people.

Include at least one person who will actually read
calldata. A quorum of executives who tap Approve because
the group chat said so is a hot wallet with extra
steps.

\section{How you operate}

Dedicated channel for signing. Backup way to reach
people. Out-of-band check before adding or replacing a
signer: video plus a signed message, not a DM.

Never copy destinations from explorer history or chat.
Poisoned lookalikes are cheap. First-time recipients get
a tiny bidirectional test. Address changes get the full
re-verify, not an edit in the address book.

Anyone can create a Safe, add your address, and park a
lookalike in your dashboard. Check the full address.
Bookmark the real one. Ignore mystery Safes.

If a key might be burned, tell the others immediately.
The remaining quorum replaces it. That is the feature.
Do not wait until you ``know for sure.''
